\section{Reproducibility Details}
\label{sec:reproducibility}

\paragraph{Models and data.}
All experiments use a frozen pretrained DiT-XL/2~\cite{peebles2022scalable} for ImageNet-derived experiments and Stable Diffusion~1.5~\cite{rombach2021highresolution} for non-ImageNet experiments, both publicly available. The complexity analysis pipeline (VAE feature extraction, K-means clustering, sigmoid mapping) uses only standard libraries (PyTorch, scikit-learn).

\paragraph{Hyperparameters.}
We report all hyperparameters: sigmoid slope $s{=}3.0$, center $c_0{=}0.6$, guidance range $[0.0, 0.06]$, complexity weights $(\alpha{=}0, \beta{=}0.5, \gamma{=}0.5, \delta{=}0)$, 50 denoising steps with a 25-step high-noise guidance window, and training with SGD (lr $0.1$, weight decay $1\text{e}{-4}$, batch size $128$). Classifier training uses 2000 epochs for 100-class experiments, ImageNette, and ImageWoof; 3000 for 10-class and CIFAR-10; 1000 for IN-1K. Data augmentation includes random resized crop, horizontal flip, color jitter, PCA lighting noise, and CutMix.

\paragraph{Random seeds.}
The main 100-class IPC=10 experiments (CAGS and Unguided on all three architectures) use 5 random seeds (seeds 0--4). All other experiments use 3 random seeds (seeds 0, 1, 2) except where noted in table captions (IN-1K ConvNet-6 uses 2 seeds; some sensitivity analysis entries use 2 seeds). Population std is used for reported standard deviations; sample std (ddof=1) is used for $t$-tests.

\paragraph{Hardware.}
All experiments run on NVIDIA RTX 4090 GPUs. The one-time complexity analysis takes approximately 22 minutes; dataset generation takes $\sim$16 minutes per configuration for 100 classes (IPC=10).

\paragraph{LLM usage.}
Large language models were used for code implementation and debugging, proof formatting, and manuscript editing. All experimental design, analysis, and scientific conclusions were conducted solely by the authors. Full disclosure is provided in the AI Use Statement in the main text.
